% (c) Nikita Lisitsa, lisyarus@gmail.com, 2026

\documentclass[10pt]{beamer}

\usepackage[T2A]{fontenc}
\usepackage[russian]{babel}
\usepackage{minted}

\usepackage[most]{tcolorbox}
\tcbuselibrary{minted}

\usepackage{graphicx}
\graphicspath{ {./images/} }

\usepackage{adjustbox}

\usepackage{color}
\usepackage{soul}
\usepackage{multicol}
\usepackage{multirow}
\usepackage{hyperref}
\usepackage{tikz}
\usetikzlibrary{arrows.meta,positioning,calc}

\usetheme{metropolis}

\definecolor{red}{rgb}{1,0,0}
\definecolor{codebg}{RGB}{29,35,49}
\setminted{fontsize=\footnotesize}

\makeatletter
\newcommand{\slideimage}[1]{
  \begin{figure}
    \begin{adjustbox}{width=\textwidth, totalheight=\textheight-2\baselineskip-2\baselineskip,keepaspectratio}
      \includegraphics{#1}
    \end{adjustbox}
  \end{figure}
}
\makeatother

\title{Язык программирования C}
\subtitle{Лекция 3: указатели на функции, виды памяти, стек, глобальные переменные, static}
\author{Лисица Никита}
\date{2026}

\setbeamertemplate{footline}[frame number]

\usemintedstyle{tango}

\begin{document}

\frame{\titlepage}

\begin{frame}[fragile]
\frametitle{``Универсальная'' сортировка}
\begin{itemize}
\item Хотим написать универсальную/обобщённую сортировку, работающую с любым типом данных: \mintinline{cpp}|int|, \mintinline{cpp}|float|, \mintinline{cpp}|const char*|, \mintinline{cpp}|struct product_t|
\pause
\item Проблемы:
\begin{itemize}
\item Как передать массив объектов произвольного типа?
\pause
\item Как переставить (\mintinline{cpp}|swap|) местами объекты произвольного типа?
\pause
\item Как сравнить (\mintinline{cpp}|compare|) объекты произвольного типа?
\end{itemize}
\end{itemize}
\end{frame}

\begin{frame}[fragile]
\frametitle{void*}
\begin{itemize}
\item \mintinline{cpp}|void| -- особый тип данных в C
\pause
\item У него нет значений (функция, возвращающая \mintinline{cpp}|void|, ничего не возвращает)
\pause
\item \mintinline{cpp}|void*| обычно используется как указатель на неизвестный тип (для него запрещена арифметика указателей)
\pause
\item Любой тип указателя приводится к \mintinline{cpp}|void*| и обратно (в C++ для обратного преобразования нужно явное приведение типа)
\pause
\item Некоторые функции (например, выделение памяти -- \mintinline{cpp}|malloc|) возвращают \mintinline{cpp}|void*|
\end{itemize}
\end{frame}

\begin{frame}[fragile]
\frametitle{void*}
\begin{itemize}
\item Мы можем использовать \mintinline{cpp}|void*|, чтобы принимать массив, который нам нужно отсортировать:
\usemintedstyle{lightbulb}
\begin{tcblisting}{listing engine=minted, minted language=cpp, minted options={fontsize=\scriptsize}, listing only, colback=codebg, colframe=codebg, rounded corners}
void sort(void *array, size_t count, ...);
\end{tcblisting}
\end{itemize}
\end{frame}

\begin{frame}[fragile]
\frametitle{Перестановка объектов (swap)}
\begin{itemize}
\item Будем считать, что перестановка сводится к перестановке байт объектов
\usemintedstyle{lightbulb}
\begin{tcblisting}{listing engine=minted, minted language=cpp, minted options={fontsize=\scriptsize}, listing only, colback=codebg, colframe=codebg, rounded corners}
void sort(void *array, size_t count, size_t elem_size, ...) {
  char *ptr = array;

  ...
  swap(array + i * elem_size, array + j * elem_size, elem_size);
  ...
}

void swap(char *p1, char *p2, size_t elem_size) {
  char *end1 = p1 + elem_size;
  while (p1 != end1) {
    ...
  }
}
\end{tcblisting}
\end{itemize}
\end{frame}

\begin{frame}[fragile]
\frametitle{Сравнение объектов (compare)}
\begin{itemize}
\item Для сравнения разных типов данных нужен разный код:
\pause
\begin{itemize}
\item Сравнение \mintinline{cpp}|int| -- оператор целочисленного сравнения
\pause
\item Сравнение \mintinline{cpp}|float| -- оператор floating-point сравнения
\pause
\item Сравнение строк (\mintinline{cpp}|const char*|) -- лексикографическое сравнение циклом (\mintinline{cpp}|strcmp|)
\pause
\item Сравнение структур \mintinline{cpp}|struct product_t| по значению поля \mintinline{cpp}|weight| -- доступ к полю + оператор сравнения
\end{itemize}
\pause
\item Как передать в функцию \mintinline{cpp}|sort| произвольный код?
\end{itemize}
\end{frame}

\begin{frame}[fragile]
\frametitle{Указатели на функции}
\begin{itemize}
\item Скомпилированный (бинарный) код программы -- тоже числа, и они хранятся где-то в памяти
\pause
\begin{itemize}
\item Гарвардская архитектура: данные и код хранятся в разных типах памяти
\pause
\item Архитектура фон Неймана: данные и код хранятся в одной и той же памяти
\end{itemize}
\pause
\item Так как это тоже память, у неё тоже есть адрес, и на неё тоже можно взять указатель
\end{itemize}
\end{frame}

\begin{frame}[fragile]
\frametitle{Указатели на функции}
\usemintedstyle{lightbulb}
\begin{tcblisting}{listing engine=minted, minted language=cpp, minted options={fontsize=\scriptsize}, listing only, colback=codebg, colframe=codebg, rounded corners}
int add3(int x) {
  return x + 3;
}

int main() {
  // Объявление указателя на функцию
  // Похоже на объявление функции, но со звёздочкой
  int (*func)(int);

  // Эти две строчки эквивалентны
  func = &add3;
  func = add3;

  // Указатель на функцию можно вызвать так же, как
  // обычную функцию
  int x = func(5);
}
\end{tcblisting}
\end{frame}

\begin{frame}[fragile]
\frametitle{Указатели на функции}
\begin{itemize}
\item На любую функцию можно взять указатель (но только на саму функцию, а не на произвольное место в коде)
\pause
\item Указатели типизированные -- знают, какие у функции аргументы и возвращаемое значение
\pause
\item Для указателей на функцию нет арифметики указателей
\end{itemize}
\end{frame}

\begin{frame}[fragile]
\frametitle{Сравнение объектов (compare)}
\usemintedstyle{lightbulb}
\begin{tcblisting}{listing engine=minted, minted language=cpp, minted options={fontsize=\scriptsize}, listing only, colback=codebg, colframe=codebg, rounded corners}
void sort(void *array, size_t count, size_t elem_size,
  int (*compare)(void*, void*)) {
  ...
}

int compare_int(void *p1, void *p2) {
  return *(int *)p1 - *(int *)p2;
}

int compare_product(void *p1, void *p2) {
  return ((product_t*)p1)->weight - ((product_t*)p2)->weight;
}


int array[100];
sort(array, 100, sizeof(int), compare_int);
\end{tcblisting}
\end{frame}

\begin{frame}[fragile]
\frametitle{Виды памяти}
\begin{itemize}
\item \textbf{Стек (stack)}
\begin{itemize}
\item Локальные переменные, параметры функций, адрес возврата
\item Код для выделения и освобождения генерирует компилятор
\item Выделяется при входе в функцию, освобождается при выходе
\end{itemize}
\pause
\item \textbf{Глобальная (статическая, static) память}
\begin{itemize}
\item Глобальные и статические переменные
\item Код для выделения и освобождения генерирует компилятор
\item Выделяется при старте программы, освобождается при её окончании
\end{itemize}
\pause
\item \textbf{Динамическая память (куча, heap)}
\begin{itemize}
\item Полностью управляется программистом
\end{itemize}
\end{itemize}
\end{frame}

\begin{frame}[fragile]
\frametitle{Адресное пространство}
\begin{itemize}
\item Точное устройство адресного пространства зависит от системы
\item Упрощённый современный линукс:
\end{itemize}
\usemintedstyle{lightbulb}
\begin{tcblisting}{listing engine=minted, minted language=cpp, minted options={fontsize=\scriptsize}, listing only, colback=codebg, colframe=codebg, rounded corners}
┌───────┬───────────┬────────────┬────────────────┬────────────────┐
│  0x0  │ 0x400000  │ 0x40xxxx   │ 0x5xxxxxxxxxxx │ 0x7ffxxxxxxxxx │
│       │           │            │                │                │
│       │   Код     │ Глобальные │  Куча (растёт  │  Стек (растёт  │
│ Пусто │ программы │ переменные │    туда ->)    │    <- туда)    │
└───────┴───────────┴────────────┴────────────────┴────────────────┘
\end{tcblisting}
\end{frame}

\begin{frame}[fragile]
\frametitle{Стек}
\begin{itemize}
\item Выделяется операционной системой при запуске прграммы (напр. 8 Мб)
\pause
\item Растёт ``вниз'' (в сторону меньших адресов)
\pause
\item Используется в основном для вызова функций
\end{itemize}
\end{frame}

\begin{frame}[fragile]
\frametitle{Стек}
\usemintedstyle{lightbulb}
\begin{tcblisting}{listing engine=minted, minted language=cpp, minted options={fontsize=\scriptsize}, listing only, colback=codebg, colframe=codebg, rounded corners}
int sum(int a, int b) {
  return a + b;
}

int main() {
  int c = 1, d = 2;
  int e = sum(c, d);
  return 0;
}
\end{tcblisting}
\pause
\usemintedstyle{tango}
\begin{itemize}
\item При вызове функции нужно
\begin{itemize}
\item Выделить место под её аргументы
\item Выделить место под её локальные переменные
\item Выделить место под её возвращаемое значение
\item Куда-то положить адрес возврата (куда должен вернуться \mintinline{cpp}|return|)
\end{itemize}
\pause
\item \(\Longrightarrow\) \textit{стек фрейм (stack frame)}
\end{itemize}
\end{frame}

\begin{frame}[fragile]
\frametitle{Стек}
\begin{itemize}
\item Стек состоит из \textit{фреймов} для всех вызванных в текущий момент функций (включая функции, вызвавшие другие функции)
\end{itemize}
\usemintedstyle{lightbulb}
\begin{tcblisting}{listing engine=minted, minted language=cpp, minted options={fontsize=\scriptsize}, listing only, colback=codebg, colframe=codebg, rounded corners}
┌────────────────────────────────────┐
│                [main]              │
│ c = 1                              │
│ d = 2                              │
│ e = (uninitialized)                │
│ return address = ??? (libc)        │
├────────────────────────────────────┤
│                [sum]               │
│ a = 1                              │
│ b = 2                              │
│ return address = 0x40143 (main+10) │
└────────────────────────────────────┘
\end{tcblisting}
\end{frame}

\begin{frame}[fragile]
\frametitle{Стек}
\begin{itemize}
\item Размер стек фрейма легко вычислить компилятору -- он знает количество аргументов и локальных переменных
\pause
\item Адрес текущей вершины стека хранится в специальном регистре SP (stack pointer)
\pause
\item Выделить N байт на стеке -- то же самое, что сделать \mintinline{cpp}|SP -= N|, т.е. \underline{очень быстро}
\pause
\item \(\Longrightarrow\) создание локальных переменных -- \underline{крайне быстрая} операция
\end{itemize}
\end{frame}

\begin{frame}[fragile]
\frametitle{Стек vs регистры}
\begin{itemize}
\item Читать/писать в стек быстро, но это всё равно работа с памятью
\pause
\item Быстрее -- передавать параметры и возвращаемое значение через регистры процессора (если они влезут)
\pause
\item Договорённость о том, какие регистры под это использовать, является частью т.н. \textit{calling convention}
\end{itemize}
\end{frame}

\begin{frame}[fragile]
\frametitle{Стек и рекурсия}
\usemintedstyle{lightbulb}
\begin{tcblisting}{listing engine=minted, minted language=cpp, minted options={fontsize=\scriptsize}, listing only, colback=codebg, colframe=codebg, rounded corners}
int factorial(int n) {
  if (n == 0) {
    return 1;
  }
  return n * factorial(n - 1);
}
\end{tcblisting}
\pause
\begin{itemize}
\item Каждый рекурсивный вызов выделяет свой фрейм
\pause
\item Если глубина рекурсии слишком большая, память стека закончится и программа упадёт с ошибкой \textit{stack overflow}
\pause
\item Почти всегда либо глубина ограничена свойствами структуры данных (напр. \(O(\log n)\) у деревьев), либо алгоритм можно переписать в виде цикла вместо рекурсии
\end{itemize}
\end{frame}

\begin{frame}[fragile]
\frametitle{Факториал без рекурсии}
\usemintedstyle{lightbulb}
\begin{tcblisting}{listing engine=minted, minted language=cpp, minted options={fontsize=\scriptsize}, listing only, colback=codebg, colframe=codebg, rounded corners}
int factorial(int n) {
  int result = 1;
  while (n > 0) {
    result *= n;
    --n;
  }
  return result;
}
\end{tcblisting}
\end{frame}

\begin{frame}[fragile]
\frametitle{Хвостовая рекурсия}
\usemintedstyle{lightbulb}
\begin{tcblisting}{listing engine=minted, minted language=cpp, minted options={fontsize=\scriptsize}, listing only, colback=codebg, colframe=codebg, rounded corners}
int factorial(int n, int result) {
  if (n == 0) {
    return result;
  }
  return factorial(n - 1, result * n);
}
\end{tcblisting}
\pause
\usemintedstyle{tango}
\begin{itemize}
\item Если рекурсивный вызов это \textbf{самая последняя операция} перед \mintinline{cpp}|return|, компилятор может переиспользовать фрейм текущей функции и не выделять новый
\pause
\item Эта оптимизация называется \textit{хвостовой рекурсией (tail call optimization)}
\end{itemize}
\end{frame}

\begin{frame}[fragile]
\frametitle{Глобальные переменные}
\begin{itemize}
\item Переменные, объявленные снаружи любой функции, называются глобальными
\pause
\item Они находятся в отдельной области памяти, и живут в течении всей жизни программы
\end{itemize}
\pause
\usemintedstyle{lightbulb}
\begin{tcblisting}{listing engine=minted, minted language=cpp, minted options={fontsize=\scriptsize}, listing only, colback=codebg, colframe=codebg, rounded corners}
int random_seed = 0;

void set_seed(int seed) {
  random_seed = seed;
}

int random() {
  random_seed = (random_seed * 13 + 113) % 43;
  return random_seed;
}
\end{tcblisting}
\end{frame}

\begin{frame}[fragile]
\frametitle{Глобальные переменные}
\begin{itemize}
\item Объявить одну и ту же глобальную переменную в нескольких файлах -- ошибка multiple definition
\pause
\item Можно исправить, пометив одно из объявлений как \mintinline{cpp}|extern|
\end{itemize}
\begin{multicols}{2}
\usemintedstyle{lightbulb}
\begin{tcblisting}{listing engine=minted, minted language=cpp, minted options={fontsize=\scriptsize}, listing only, colback=codebg, colframe=codebg, rounded corners}
// a.cpp

// Аналог определения
// (definition) функции
int global = 0;
\end{tcblisting}
\columnbreak
\begin{tcblisting}{listing engine=minted, minted language=cpp, minted options={fontsize=\scriptsize}, listing only, colback=codebg, colframe=codebg, rounded corners}
// b.cpp

// Аналог объявления
// (declaration) функции
extern int global;
\end{tcblisting}
\end{multicols}
\end{frame}

\begin{frame}[fragile]
\frametitle{Глобальные переменные + static}
\begin{itemize}
\item Глобальная переменная, объявленная со словом \mintinline{cpp}|static|, видна \textbf{только в этом файле} и не приводит к ошибке линковки
\end{itemize}
\begin{multicols}{2}
\usemintedstyle{lightbulb}
\begin{tcblisting}{listing engine=minted, minted language=cpp, minted options={fontsize=\scriptsize}, listing only, colback=codebg, colframe=codebg, rounded corners}
// a.cpp



static int global = 0;
\end{tcblisting}
\columnbreak
\begin{tcblisting}{listing engine=minted, minted language=cpp, minted options={fontsize=\scriptsize}, listing only, colback=codebg, colframe=codebg, rounded corners}
// b.cpp

// Ошибки нет: в каждом файле
// своя переменная global
static int global = 0;
\end{tcblisting}
\end{multicols}
\end{frame}

\begin{frame}[fragile]
\frametitle{Статические переменные}
\begin{itemize}
\item Переменные внутри функции, объявленные со словом \mintinline{cpp}|static|
\item Живут вместе с глобальными переменными, но доступны только внутри этой функции
\end{itemize}
\usemintedstyle{lightbulb}
\begin{tcblisting}{listing engine=minted, minted language=cpp, minted options={fontsize=\scriptsize}, listing only, colback=codebg, colframe=codebg, rounded corners}
void foo() {
  // Инициализируется один раз
  // Значение сохраняется даже после выхода из функции
  static int call_count = 0;

  printf("I was called %i times\n", call_count);
  ++call_count;
}
\end{tcblisting}
\end{frame}

\begin{frame}[fragile]
\frametitle{Глобальные переменные}
\begin{itemize}
\item При компиляции компилятор может легко посчитать, сколько памяти нужно под все глобальные/статические переменные в этом файле
\pause
\item При линковке линкофщик может легко просуммировать необходимую память под все глобальные/статические переменные всех файлов
\pause
\item Эта память выделится в начале работы программы и освободится при её завершении
\end{itemize}
\end{frame}

\begin{frame}[fragile]
\frametitle{Глобальные переменные: недостатки}
\begin{itemize}
\item Потенциальный конфликт имён, если два программиста завели в разных частях программы одинаковые глобальные переменные
\pause
\item Трудно анализировать, какие части программы что используют и на что могут повлиять
\pause
\item Порядок инициализации глобальных переменных в рамках всей программы \textbf{не определён} (только в рамках одного файла)
\pause
\item Считаются плохим тоном/антипаттерном, рекомендуется их не использовать без крайней необходимости
\end{itemize}
\end{frame}

\end{document}
